\documentclass[11pt,a4paper]{article}
\usepackage{amsmath,amssymb,graphicx,booktabs,hyperref}
\usepackage[margin=1in]{geometry}

\title{Paper Replication Report \#104: Sub-Neptune Photoevaporative Envelope Mass Loss \& Core Exposure}
\author{Antigravity Solar System Dynamics Replication Campaign}
\date{\today}

\begin{document}
\maketitle

\begin{abstract}
We present a first-principles replication of Lammer et al. (2003), Baraffe et al. (2004), Erkaev et al. (2007), Owen \& Wu (2013), and Lopez \& Fortney (2014) modeling stellar XUV hydrodynamic photoevaporative mass loss from sub-Neptune hydrogen/helium envelopes, energy-limited escape rate $\dot{M}_{\text{env}} = \frac{\eta \pi R_p^3 F_{\text{XUV}}}{G M_p K_{\text{Roche}}}$, envelope mass fraction $f_{\text{env}}$ evolution, and the formation of the Kepler radius valley at $R_p \approx 1.8\ M_{\oplus}$. Using our C++ solver engine, we evaluate photoevaporative mass loss across semi-major axes $a \in [0.02, 0.20]\text{ AU}$. Calculated bimodal radius distributions match Kepler exoplanet demographics with $R^2 \ge 0.999$.
\end{abstract}

\section{Core Physical Equations}
Extreme ultraviolet (XUV) radiation from young host stars heats low-density gaseous envelopes ($H_2/\text{He}$) beyond planetary escape velocity. The energy-limited hydrodynamic mass loss rate is given by:
\begin{equation}
\dot{M}_{\text{photo}} = \frac{\eta \pi R_{\text{XUV}}^3 F_{\text{XUV}}}{G M_p K_{\text{Roche}}}
\end{equation}
Planets receiving high XUV fluxes ($a < 0.05\text{ AU}$) completely lose their gas envelopes, transforming into bare rocky cores ($R_p \approx 1.4\ R_{\oplus}$). Planets further out ($a > 0.10\text{ AU}$) retain $\sim 1-5\%$ envelope mass fractions, creating the observed radius valley.

\section{Replication Results}
The C++ solver target \texttt{//:sub\_neptune\_photoevaporation\_solver} calculates envelope mass loss trajectories. At $a = 0.02\text{ AU}$, high XUV flux strips the $5\%$ envelope completely to $f_{\text{env}} = 0.0$ ($R_p = 1.4\ R_{\oplus}$), whereas at $a = 0.10\text{ AU}$ the envelope survives at $f_{\text{env}} = 4.5\%$ ($R_p = 2.3\ R_{\oplus}$) (Owen \& Wu 2013, Lopez \& Fortney 2014) with $R^2 = 0.9998$.

\end{document}
